\chapter{Comparing Proofs and Foundations}\label{ch:proof-comparison}

\begin{center}
\href{proof-bench/index.html}{\textbf{Open the proof-size laboratory}}
\end{center}

A short final theorem can be an application of a large pre-existing library.
Conversely, a reusable general theorem should not be charged from scratch at
every application. The comparison therefore fixes one computational statement
and measures several named derivations, rather than equating source length,
import count, and mathematical simplicity.

\section{Contracts before measurements}
Every case records its exact computation, rational domain, and hypotheses.
The exporter checks that the complete Lean theorem types of its alternatives
are definitionally equal. It rejects a proof that calls its alternative, a
Mathlib dependency in a route marked native, and transitive admitted proofs.
These checks do not replace reviewing the mathematical contract: an extra
assumption may make a statement uninteresting even if every proof has it.

The initial registry contains the three cosine endpoint proofs, the three
validity proofs of that same quadrature program, and two proofs of the sine
value at zero. This is \emph{one calculus example and one calibration}, not a
corpus-wide finding about the simplest calculus foundation.

\section{Three explicit accounting baselines}
Let $D(p)$ contain the declaration $p$ and the transitive references in its
stored type and body. Let $T$ be the union of statement-prerequisite closures
for the matched alternatives, and $H$ the intersection of their full closures.
We display the full cost $D(p)$, the cost beyond the stated objects
$D(p)\setminus T$, and the route-specific cost $D(p)\setminus H$.
The removed baseline is declared explicitly; it has not become mathematically
free. Display-only companion declarations are not measurement roots.

Each set is partitioned into the native project, comparison bridges, Mathlib,
Lean/Std, and other dependencies. We report stored expression-tree occurrences
and distinct subexpressions within each body, with type-expression sizes
separate. Named constants are not unfolded; duplicate references to the same
declaration do not charge its body twice. Native-computation certificates and
numeric literals remain atomic references, so trust reports accompany sizes.
Available source-range unions are a formatting-dependent diagnostic with
explicit coverage, not a substitute for expression measurements.

\section{Reuse over a corpus}
For a route applied to ordered cases $p_1,\ldots,p_k$, the accumulated library
is $U_k=\bigcup_{i\le k}D(p_i)$. Its incremental cost at case $k$ is
$U_k\setminus U_{k-1}$. The final union is order-independent; marginal costs
are not. The first live example accumulates cosine identity and validity,
without counting them as two unrelated calculus examples.

The predeclared next families are polynomial FTC uniformly in degree,
arctangent-kernel integration on rational subintervals, the actual sine
secant derivative, square-root inversion and differentiation away from zero,
global rational-angle addition identities, exponential integration, computable
complex exponential, and linear ODE solution/uniqueness. Their statuses remain
planned until paired proofs and bridges pass the same gates. Varying a degree
or endpoint must not manufacture a collection of apparently independent wins.

\section{What a smaller footprint would establish}
A smaller used formal library for a fixed contract is evidence of a narrower
dependency footprint, not a proof of conceptual minimality. The native route
uses explicit rational programs and finite invariants; Mathlib supplies much
broader real analysis. The computational target representation also places
bridge work on the Mathlib route, which is why bridge cost is separate.

Both use the same Lean kernel. Avoiding the completed real type does not imply
absence of ordinary Lean axioms or native-computation assumptions. Nor does
reaching Mathlib's real type count how many mathematical completeness theorems
are used. Compare domains, regularity hypotheses, effectiveness, theorem reuse,
trust, and runtime alongside code size. Kernel timing, elaboration timing, and
numerical execution timing are separate future measurements, not inferred
from the expression counts on this page.
